\documentclass[11pt]{article}
\usepackage[margin=1in]{geometry}
\usepackage{amsmath,amssymb,amsthm}
\usepackage[hidelinks]{hyperref}

\newtheorem{theorem}{Theorem}

\newcommand{\PP}{\mathbb P}
\newcommand{\RR}{\mathbb R}

\begin{document}

\begin{center}
{\Large A smaller-threshold tail consequence for Bernoulli sections}\\[0.5em]
{\small Run \texttt{fa\_banach\_001}, arXiv:math/0601369, literature-implied partial subcase}
\end{center}

\section*{Source Question}

Artstein-Avidan, Friedland, Milman and Sodin study a \(p\times n\)
Bernoulli sign matrix \(X\), with \(p=(1-\delta)n\), and
\[
  S=n^{-1}XX^T .
\]
Their Theorem 2 proves
\[
  \PP\{\lambda_{\min}(S)\leq \delta^2/8\}
  \leq \exp(-c n^{1/6}\delta)
\]
under the stated lower bound on \(\delta\).  Immediately after comparing
this with the Litvak--Pajor--Rudelson--Tomczak-Jaegermann estimate, they ask
whether this left-hand side is in fact as small as
\[
  \exp(-n\delta^C/C)
\]
for some universal \(C>0\) \cite{AFMS}.

\section*{Supporting Literature}

Dabagia and Fernandez prove a rectangular smallest-singular-value lower-tail
theorem for random matrices with independent entries and bounded
\(2+\beta\) moments.  In particular, for an \(N\times m\) matrix \(A\) in
their setting, \(N>m\), their theorem gives constants \(C,c>0\) such that
\[
  \PP\left\{\sigma_m(A)\leq
    \varepsilon(\sqrt{N+1}-\sqrt m)\right\}
  \leq (C\varepsilon)^{N-m+1}+e^{-cN}
\]
for all \(\varepsilon>0\) \cite{DF}.  Bernoulli sign entries satisfy the
moment hypotheses with universal constants.

\section*{Implied Partial Answer}

\begin{theorem}
There are universal constants \(c_0,c_1,c_2>0\) such that, for every
\(0<\delta<1/2\), \(p=(1-\delta)n\), and \(p\times n\) Bernoulli sign matrix
\(X\),
\[
  \PP\left\{\lambda_{\min}\left(n^{-1}XX^T\right)\leq c_0\delta^2\right\}
  \leq e^{-c_1\delta n}+e^{-c_2 n}.
\]
\end{theorem}

\begin{proof}
Apply the Dabagia--Fernandez theorem to \(A=X^T\), viewed as an \(n\times p\)
matrix.  Choose a fixed \(\varepsilon_0>0\) so small that
\((C\varepsilon_0)^{n-p+1}\leq e^{-c_1(n-p+1)}\).  Then
\[
  \PP\{\sigma_p(X)\leq
  \varepsilon_0(\sqrt{n+1}-\sqrt p)\}
  \leq e^{-c_1(n-p+1)}+e^{-c_2n}.
\]
Since \(p=(1-\delta)n\) and \(0<\delta<1/2\),
\[
  \sqrt{n+1}-\sqrt p\geq \sqrt n-\sqrt{(1-\delta)n}
  = \sqrt n(1-\sqrt{1-\delta})\geq \frac{\delta\sqrt n}{2}.
\]
Thus the event
\[
  \lambda_{\min}(n^{-1}XX^T)=\sigma_p(X)^2/n
  \leq \varepsilon_0^2\delta^2/4
\]
is contained in the event estimated above.  Set
\(c_0=\varepsilon_0^2/4\).
\end{proof}

\section*{Status and Limitations}

This is a literature-implied partial subcase.  It gives the exponential tail
shape requested in the 2006 question only after replacing the source threshold
\(\delta^2/8\) by a smaller universal multiple \(c_0\delta^2\).  The exact
left-hand side in the source question, with threshold \(\delta^2/8\), is not
claimed to be answered here.

The supporting paper does not explicitly cite arXiv:math/0601369 as the
question being answered; the implication is obtained by identifying the
source's covariance eigenvalue with the square of the smallest singular value
of the transposed rectangular Bernoulli matrix.

\begin{thebibliography}{9}

\bibitem{AFMS}
S. Artstein-Avidan, O. Friedland, V. Milman, and S. Sodin,
\emph{Polynomial bounds for large Bernoulli sections of \(\ell_1^N\)},
arXiv:math/0601369, 2006.

\bibitem{DF}
M. Dabagia and M. Fernandez,
\emph{The smallest singular value of inhomogenous random rectangular matrices},
arXiv:2408.14389, version dated 2026.

\end{thebibliography}

\end{document}

